\label{sec:interdep}
The question was how an optimal system could fail to couple the sampled weights and its estimate, perhaps dynamically. There are four
exact couplings, ordered by how much they are forced.

\subsection{Forced: the frame}
Section~\ref{sec:growth} shows that the representation must be re-centred at every layer by the instance's own law. The frame of layer $l$ is a
function of $W_1,\dots,W_l$, and it changes as the computation proceeds.

\subsection{Free: the estimator learns its own correction from its own rows}
\begin{theorem}[row-exchangeable self-calibration; proved]\label{thm:selfcal}
Fix a layer $l$ and the upstream state $U$, meaning the weights below $l$ and the estimator's state up to layer $l-1$. For each row $w_i$ of $W_l$,
let $\pi_i=\pi(w_i;U)\in\R^p$ be a cheap per-row feature vector and $\beta_i=\beta(w_i;U)$ an expensive per-row quantity, for instance a higher-order birth.
Choose a uniformly random set $S$ of $q$ rows, fit $\hat a$ by least squares of $\beta$ on $\pi$ over $S$, and set $\hat\beta_i=\hat a\cdot\pi_i$ off $S$. Then:
\begin{enumerate}[leftmargin=*]
\item Given $U$, the pairs $(\pi_i,\beta_i)$ are i.i.d., because the rows are.
\item $\hat a$ converges to the population regression $a_*$, with excess error $O(p/q)$ relative to the residual.
\item The estimator's added mean squared error at layer $l$ is $(1-\frac qn)(1-R^2)\,\E\beta^2\,(1+O(p/q))$, where $R^2$ is the population fit.
\item By the error cocycle, this error reaches the output additively, multiplied by the layer's downstream gain.
\end{enumerate}
\end{theorem}
\begin{proof}
(1) The rows are i.i.d.\ and independent of $U$. (2) and (3) are least squares on i.i.d.\ samples. (4) is Theorem~3.2 of stage~15.
\end{proof}
In this precise sense an estimator learns its own correction from its own weights. The training set is the instance's rows, and what makes
it valid is exchangeability, not any independence of weights and estimate. Corollary~3.3 of stage~15 bounds $R^2$ by the isotropic share
whenever $\pi$ holds only per-neuron or field features. A large $R^2$ needs anisotropic proxies, such as a short-window or first-order
version of $\beta$ itself.

\emph{Where it should pay, and the test.} It pays where per-row costs dominate and such a proxy exists. Late-layer covariance births are
the case: layers $11$--$16$ hold $81\%$ of the exact chain's births, and $59\%$ of the last-layer birth is covariance.
\begin{itemize}[leftmargin=*]
\item The expensive quantity: the second-order covariance-birth readout of a row (the $(2,2)$, $(3,1)$ and $\kappa_3$-pair diagrams),
computed on $q\approx n/8$ rows.
\item The proxy: the first-order readout plus the field jets.
\item The prediction: the share $R^2$ of the channel is recovered at about an eighth of its readout cost.
\end{itemize}

\subsection{Free: geometry coupled to the readout}
The holonomy note's principle $\mathcal V(w)=\langle f,\mathsf K_w^{-1}f\rangle$ is convex in the weights $w$, and its optimality conditions equalize the edge
energies of the Poisson solution. Our analogue is stage~15's allocation:
\begin{itemize}[leftmargin=*]
\item The predicted MSE is convex in the capture fractions.
\item At the optimum the marginal MSE reduction per FLOP is equal across the (layer, row, diagram) choices used.
\item The Poisson solution is the adjoint: the downstream gains.
\end{itemize}
The two are coupled self-consistently. The allocation depends on the adjoint, and the adjoint depends on what was computed. Two passes,
or one pass with revision, solve it.

\subsection{Free: task-preserving instruments}
The companion's condition $\sum_\alpha\mathcal I_\alpha^*(A_\alpha)=A$ has a classical realization.
\begin{itemize}[leftmargin=*]
\item Branch on the sign of a chosen functional $u\cdot x$. By symmetry each branch has mass $\frac12$.
\item Run each branch with its own Wick frame. Then $\E f=\sum_\alpha P(\alpha)\E[f\mid\alpha]$ holds exactly.
\item A branch resolves all orders of non-Gaussianity along $u$, at twice the cost.
\end{itemize}
Its value depends on concentration, and the residual is not concentrated. In the experimental session, making the covariance exact in the top
$128$ collective modes gave only $16$--$32\%$, which a branching tree could approach only with $2^{128}$ leaves. Instruments are the right tool for
a coherent residual and the wrong one for an incoherent residual.

\subsection{The design space}
An estimator is a point in (frame family) $\times$ (partition blocks evaluated on the instance) $\times$ (rows) $\times$ (layers) $\times$ (instruments).
\begin{itemize}[leftmargin=*]
\item \emph{The standard chain.} Gaussian centred frames, pair and three-block diagrams, all rows, all layers, no instruments.
\item \emph{Forced by this stage.} Co-evolving frames and transport memory.
\item \emph{Ruled out here.} Leaf estimators with global energies, and minimal-sensitivity Gaussian frames.
\item \emph{Open, with protocols.} Non-Gaussian frame families, diagram selection by annealed two-copy variances,
self-calibration over rows, and late-layer allocation.
\end{itemize}
